% Rigid methane (closing pages, unnumbered): three panels, titles only, no leader labels.
% Upper left: methane, CH4, the reference X0 in its body frame, drawn small with the suite's molecule primitive seen
% with the ball's own camera (mol3d-ch4-ball), so that its bond 1 points the way the lift arrow in the ball does;
% the C2 axes x, y, z through the midpoints of opposite edges as thin lines behind it; its title in the corner.
% Middle, large, close to the molecule since the two share one orientation (author) (make_data: CAM_BALL, azimuth 108, elevation 21, chosen by scan to sit 27
% degrees from every symmetry axis of the cube so that nothing overlaps, bond 1 toward the viewer): the orientation
% space SO(3) as the axis-angle ball, direction the axis and distance the angle, radius pi drawn 2.4 cm; the sphere
% teal with its three great circles (the skin is the half-turns, each point its own antipode). The twelve rotations
% of T as red dots: the centre, the eight third-turns at 2pi/3 along the body diagonals, the three half-turns on the
% skin, each drawn once at its far representative and pale, so that it reads as distant (author). The cell of the
% centre as the schematic exact in Rodrigues coordinates: the gold wire octahedron dual to the thin ink cube spanned
% by the third-turns, its vertices (the quarter-turns, ink rings) on the cube's face centres; in the ball the true
% edges bow outward (the caption says so). The lift of the third-turn loop as the one heavy arrow from the centre
% through a face to the dot beyond: its inner part drawn before the gold edges so that the edges in front of it
% cover it, its outer part after. Solid for what faces the viewer, dashed for what lies behind.
% Right, with room: the covering in figure 11's grammar: the gold base patch with the loop l_g through its point (a
% line, no arrowhead, as on the cover), three sheets r, r R_g, r R_g^2 of the twelve, the lift as the same curve raised from the
% dot on sheet r to the dot on sheet r R_g (the same curve, so that it projects down onto the loop); one vertical
% line of projection through every dot, hidden behind each sheet and re-emerging at its dot, running on past the top
% sheet with a vertical ellipsis (author): the stack goes on. The sheets are the cartoon of figure 11.
% The two bottom texts are left-aligned on one line with a gap between them (author).
% Every coordinate from figures/data (make_data); the geometry recomputed in checks/verify_methane.py.
\documentclass[10pt,border=1mm]{standalone}
\usepackage[T1]{fontenc}
\usepackage{figurestyle}
\input{primitives.tex}
\input{molecules.tex}
\input{numbers.tex}
\newcommand\csheet[5]{\path[fill=#5] ({#1-#3},#2)--({#1+#3},#2)--({#1+#3+.6},{#2+.6})--({#1-#3+.6},{#2+.6})--cycle; \draw[#4] ({#1-#3},#2)--({#1+#3},#2)--({#1+#3+.6},{#2+.6})--({#1-#3+.6},{#2+.6})--cycle;}
\begin{document}
\begin{tikzpicture}[plate,
  thin/.style={line width=.3pt,draw=PlateInk},
  thindash/.style={line width=.3pt,draw=PlateInk,dash pattern=on 1.5pt off 1.1pt},
  teal/.style={line width=.45pt,draw=ColBandLine},
  tealthin/.style={line width=.35pt,draw=ColBandLine},
  tealdash/.style={line width=.35pt,draw=ColBandLine,dash pattern=on 1.5pt off 1.1pt},
  gold/.style={line width=.65pt,draw=ColGoldStroke},
  golddash/.style={line width=.55pt,draw=ColGoldStroke,dash pattern=on 1.8pt off 1.2pt},
  blurb/.style={sub,anchor=north west,align=left}]
\path[use as bounding box] (0,0) rectangle (15,8.8);
% ================================================ upper left: methane in its body frame, seen as the ball is seen
\node[blurb] at (.15,8.7) {Methane, $\mathrm{CH_4}$: the reference $\X_0$ in its body frame,\\with the $C_2$ axes $x,y,z$ through the midpoints of opposite edges.};
\begin{scope}[shift={(1.9,6.4)}]
  \foreach \a/\lab/\anc in {X/$x$/south,Y/$y$/west,Z/$z$/west} {
    \draw[thin] ({-1.4*\csname baAxis\a x\endcsname},{-1.4*\csname baAxis\a y\endcsname})--({1.4*\csname baAxis\a x\endcsname},{1.4*\csname baAxis\a y\endcsname});
    \node[sub,anchor=\anc,inner sep=1.5pt] at ({1.47*\csname baAxis\a x\endcsname},{1.47*\csname baAxis\a y\endcsname}) {\lab};}
  \begin{scope}\molsolid\mollabelsinside\pic[scale=.8] at (0,0) {mol3d-ch4-ball};\end{scope}
\end{scope}
% ================================================ middle: the orientation ball, large
\begin{scope}[shift={(5.9,5.05)},scale=1.1]
  \foreach \pl in {XY,YZ,ZX} {\draw[tealdash] plot[smooth] coordinates {\csname ballArc\pl Back\endcsname};}
  \foreach \p/\q in \ballCubeHidden {\draw[thindash] (\csname \p x\endcsname,\csname \p y\endcsname)--(\csname \q x\endcsname,\csname \q y\endcsname);}
  \foreach \p/\q in \ballCellHidden {\draw[golddash] (\csname \p x\endcsname,\csname \p y\endcsname)--(\csname \q x\endcsname,\csname \q y\endcsname);}
  \draw[line width=.9pt,draw=PlateInk] (0,0)--(\ballFax,\ballFay);
  \foreach \p/\q in \ballCellVisible {\draw[gold] (\csname \p x\endcsname,\csname \p y\endcsname)--(\csname \q x\endcsname,\csname \q y\endcsname);}
  \draw[tpath,shorten >=2.5pt] (\ballFax,\ballFay)--(\ballCax,\ballCay);
  \foreach \p/\q in \ballCubeVisible {\draw[thin] (\csname \p x\endcsname,\csname \p y\endcsname)--(\csname \q x\endcsname,\csname \q y\endcsname);}
  \foreach \pl in {XY,YZ,ZX} {\draw[tealthin] plot[smooth] coordinates {\csname ballArc\pl Front\endcsname};}
  \draw[teal] (0,0) circle[radius=\ballR];
  \foreach \k in {a,b,c,d} {\fill[ColDot] (\csname ballC\k x\endcsname,\csname ballC\k y\endcsname) circle[radius=.06]; \fill[ColDot] (\csname ballD\k x\endcsname,\csname ballD\k y\endcsname) circle[radius=.06];}
  \foreach \s in \ballSkinBack {\fill[ColDot!45] (\csname ballS\s x\endcsname,\csname ballS\s y\endcsname) circle[radius=.06];}
  \fill[ColDot] (0,0) circle[radius=.06];
  \foreach \a in {x,y,z} {\foreach \t in {p,m} {\draw[fine,fill=white] (\csname ballV\a\t x\endcsname,\csname ballV\a\t y\endcsname) circle[radius=.04];}}
\end{scope}
\node[blurb] at (3.0,2.3) {The orientation space $SO(3)$ as a ball,\\direction the axis and distance the angle:\\the twelve rotations of $T$,\\the cell $\mathcal U$ at $q=0$ in gold, and the lift\\$r\to r\Rmat_g$ of the third-turn loop.};
% ================================================ right: the covering, with room
\begin{scope}[shift={(11.7,3.3)}]
  \coordinate (bx) at (.5,.3);
  \csheet{0}{0}{1.2}{heavy}{ColRetained!50}
  \draw[tlift] (bx) .. controls (-.7,.85) and (-.7,-.25) .. (bx);
  \node[sub,anchor=north,inner sep=1.5pt] at (-.15,-.08) {$\ell_g$};
  % one vertical line of projection through every dot: drawn under the sheets, so that each sheet hides it below
  % its dot, and redrawn above each dot, so that it re-emerges; it runs on past the top sheet
  \draw[hidden,shorten <=2.3pt] (bx)--(.5,5.6);
  \foreach \k/\lab in {0/{r},1/{r\Rmat_g},2/{r\Rmat_g^{\,2}}} {
    \pgfmathsetmacro{\y}{1.4+1.3*\k}
    \csheet{0}{\y}{1.2}{fine}{white}
    \coordinate (p\k) at (.5,{\y+.3});
    \draw[hidden,shorten <=2.3pt] (p\k)--(.5,{\y+1.3});
    \node[sub,anchor=west,inner sep=1pt] at (1.9,{\y+.3}) {$\lab$};}
  \node[sub,anchor=south,inner sep=0pt] at (.5,5.75) {$\vdots$};
  \draw[tlift] (p0) .. controls (-.95,{1.4+.3+.75}) and (-.95,{2.7+.3-.35}) .. (p1);
  \foreach \k in {0,1,2} {\fill[ColDot] (p\k) circle[radius=.06];}
  \fill[ColDot] (bx) circle[radius=.06];
\end{scope}
\node[blurb] at (9.0,2.3) {The covering $SO(3)\to SO(3)/T$,\\as in Figure~11: three of the twelve\\sheets, one per rotation in $T$; the loop $\ell_g$\\in the base, and its lift from sheet $r$ to $r\Rmat_g$.\\$\mathrm A$ sees $1$, ${}^1\mathrm E$ sees $\omega$, ${}^2\mathrm E$ sees $\omega^2$,\\and $\mathrm T$ cycles $x,y,z$.};
\end{tikzpicture}
\end{document}
